\documentclass[UTF8,a4paper,12pt]{ctexrep}

\usepackage[margin=2.3cm]{geometry}
\usepackage{xcolor}
\usepackage{hyperref}
\usepackage{booktabs}
\usepackage{longtable}
\usepackage{tabularx}
\usepackage{array}
\usepackage{enumitem}
\usepackage{listings}
\usepackage{amsmath}
\usepackage{amssymb}
\usepackage{tikz}
\usetikzlibrary{arrows.meta,positioning,fit,calc}

\hypersetup{
  colorlinks=true,
  linkcolor=blue!60!black,
  urlcolor=blue!60!black,
  citecolor=blue!60!black,
  pdftitle={四篇室内定位与多传感器融合论文总结},
  pdfauthor={RSSI/IMU Go2 Project}
}

\definecolor{codebg}{RGB}{248,248,248}
\definecolor{notebg}{RGB}{245,249,255}
\definecolor{warnbg}{RGB}{255,248,238}
\definecolor{goodgreen}{RGB}{28,120,70}
\definecolor{deepblue}{RGB}{30,88,170}
\definecolor{deeporange}{RGB}{190,96,20}

\lstdefinestyle{shellstyle}{
  basicstyle=\ttfamily\small,
  backgroundcolor=\color{codebg},
  frame=single,
  framerule=0.2pt,
  rulecolor=\color{black!25},
  breaklines=true,
  columns=fullflexible,
  keepspaces=true,
  showstringspaces=false
}
\lstset{style=shellstyle}

\newcommand{\code}[1]{\texttt{\detokenize{#1}}}
\newcommand{\file}[1]{\texttt{\detokenize{#1}}}
\newcommand{\topic}[1]{\texttt{\detokenize{#1}}}
\newcommand{\pkg}[1]{\texttt{\detokenize{#1}}}
\newcommand{\paperid}[1]{\texttt{\detokenize{#1}}}
\newcolumntype{Y}{>{\raggedright\arraybackslash}X}
\newcolumntype{L}[1]{>{\raggedright\arraybackslash}p{#1}}

\newenvironment{notebox}[1]
{\begin{center}\begin{minipage}{0.94\linewidth}\color{deepblue}\hrule\vspace{0.6em}\textbf{#1}\par\vspace{0.4em}\color{black}}
{\vspace{0.4em}\hrule\end{minipage}\end{center}}

\newenvironment{warnbox}[1]
{\begin{center}\begin{minipage}{0.94\linewidth}\color{deeporange}\hrule\vspace{0.6em}\textbf{#1}\par\vspace{0.4em}\color{black}}
{\vspace{0.4em}\hrule\end{minipage}\end{center}}

\setlist[itemize]{leftmargin=2em,itemsep=0.25em,topsep=0.35em}
\setlist[enumerate]{leftmargin=2em,itemsep=0.25em,topsep=0.35em}
\renewcommand{\arraystretch}{1.22}
\setlength{\tabcolsep}{4pt}

\title{四篇室内定位与多传感器融合论文总结\\
\large 面向 Go2 RSSI/IMU/Nav2 项目的背景、模型、算法与数据流}
\author{RSSI/IMU Go2 Project Notes}
\date{2026 年 6 月 16 日}

\begin{document}
\maketitle
\tableofcontents

\chapter{论文一：HGP-RL 分布式分层高斯过程 WiFi 相对定位}

\section{基本信息}

\begin{itemize}
  \item 论文：\emph{HGP-RL: Distributed Hierarchical Gaussian Processes for Wi-Fi-based Relative Localization in Multi-Robot Systems}
  \item 文件：\file{code/paper/2307.10614v2.pdf}
  \item 主题：WiFi RSSI、多机器人相对定位、Gaussian Process Regression、分层推理、AP-oriented 坐标变换。
\end{itemize}

\section{研究背景}

多机器人系统在物流、巡检、救援、编队、覆盖和 rendezvous 等任务中，需要知道其他机器人相对自己的位置。传统方案常依赖相机、RGB-D、LiDAR 或多机器人 SLAM，但这些方案有几个问题：

\begin{itemize}
  \item 传感器贵、重、耗电，算力需求高。
  \item 多机器人 SLAM 常需要地图重叠、特征匹配和 loop closure。
  \item 在 GPS-denied 或希望保护全局位置隐私的环境里，绝对位置并不总是可用。
  \item 对很多协作控制任务来说，相对定位足够，不一定需要完整全局地图。
\end{itemize}

WiFi RSSI 是一个诱人的替代信息源：几乎所有机器人或嵌入式平台都可获得 WiFi 信号，RSSI 与距离存在粗略相关性。但 RSSI 的困难也很明显：室内多径、遮挡、非视距、背景噪声会让 RSSI 场出现局部假峰和强波动。因此作者没有直接用 path-loss 公式，而是用高斯过程回归去学习 RSSI 空间场，并进一步设计分层推理降低实时搜索成本。

\section{问题定义}

假设有 \(R\) 个机器人都连接到同一个固定 WiFi AP。第 \(i\) 个机器人只能在自己的局部坐标系中表示位置，记为 \({}^{i}p_i\)。目标是让机器人 \(i\) 估计任意邻居机器人 \(j\) 在自己坐标系中的相对位置 \({}^{i}p_j\)。

机器人可以共享的信息只有：

\begin{itemize}
  \item 自己局部坐标系里的当前位置，来自 odometry/IMU；
  \item 自己估计到的 AP 位置；
  \item 这些信息仍然都在各自局部坐标系中表达。
\end{itemize}

这意味着共享数据不能直接相减，需要经过坐标变换。

\section{模型：RSSI 的 Gaussian Process Regression}

论文将位置到 RSSI 的映射看成未知函数：

\[
f: p=(x,y) \mapsto \mathrm{RSSI}.
\]

已采集的数据集为：

\[
\mathcal{D}=\{(x_q,y_q)\}_{q=1}^{M},
\]

其中 \(x_q\) 是机器人采样位置，\(y_q\) 是该位置测得的 RSSI。GP 由均值函数和核函数定义。论文采用常数均值：

\[
m(x)=\mu,
\]

核函数采用平方指数核，也就是 RBF 核：

\[
k(x,x')=\sigma_f^2 \exp\left(-\frac{\|x-x'\|^2}{2l^2}\right).
\]

这个核的直觉是：空间上越近的点，RSSI 越相关；越远的点，相关性越弱。对室内 RSSI 来说这不是完美真实，但作为平滑先验足够实用。

GP 推理输出两个量：

\begin{itemize}
  \item \(\mu_M(q_\star)\)：测试位置 \(q_\star\) 的预测 RSSI 均值；
  \item \(\sigma_M^2(q_\star)\)：预测不确定性。
\end{itemize}

训练复杂度约为 \(O(M^3)\)，单点推理复杂度约为 \(O(M^2)\)。因此如果在一个大区域内对大量高分辨率网格点逐点推理，代价会很高。

\section{算法一：分层 AP 搜索}

如果用普通 dense search，AP 位置可以写成：

\[
{}^{i}p_{\mathrm{AP}}^\star=\arg\max_{p\in {}^{i}B}\mu_M(p).
\]

也就是在机器人 \(i\) 的搜索边界 \({}^{i}B\) 内，找 GP 预测 RSSI 均值最大的点作为 AP。问题是若每个维度有 \(N\) 个网格点，二维空间就有 \(N^2\) 个候选点。搜索越精细，推理成本越高。

HGP-RL 的关键是 coarse-to-fine：

\begin{enumerate}
  \item 第一层用粗分辨率在较大范围内找 AP 的大致位置。
  \item 第二层以上一层结果为中心，结合上一层不确定性，缩小边界并提高分辨率。
  \item 第三层、第四层继续细化。
  \item 最终位置由各层偏移累加得到。
\end{enumerate}

复杂度从全分辨率的 \(O(N^d)\) 量级下降到：

\[
O(K(\lambda N)^d),
\]

其中 \(K\) 是层数，\(\lambda N \ll N\) 是每层稀疏后的网格规模。论文实验中使用 4 层分辨率：

\[
[0.1\text{ m}, 0.05\text{ m}, 0.025\text{ m}, 0.0125\text{ m}].
\]

论文还比较了梯度搜索。梯度搜索在真实 RSSI map 中容易陷入局部最大值，尤其多径和反射会造成假峰。分层搜索从粗到细先锁定全局大区域，再局部细化，因此更不容易被局部噪声峰吸走。

\section{算法二：AP-oriented 相对定位}

所有机器人都在定位同一个物理 AP。设机器人 \(i\) 估计的 AP 位置为 \({}^{i}p_{\mathrm{AP}}\)，机器人 \(j\) 估计的 AP 位置为 \({}^{j}p_{\mathrm{AP}}\)。如果知道从 \(j\) 坐标系到 \(i\) 坐标系的旋转矩阵 \({}^{i}R_j\)，则同一条“机器人到 AP”的向量在两个坐标系中满足：

\[
({}^{i}p_j-{}^{i}p_{\mathrm{AP}})
=
{}^{i}R_j({}^{j}p_j-{}^{j}p_{\mathrm{AP}}).
\]

因此机器人 \(i\) 可以求出机器人 \(j\) 在自己坐标系中的位置：

\[
{}^{i}p_j
=
{}^{i}p_{\mathrm{AP}}
+{}^{i}R_j({}^{j}p_j-{}^{j}p_{\mathrm{AP}}).
\]

这一公式很重要。它说明论文并不需要全局地图，也不需要机器人之间直接测距。AP 成了一个共同锚点，机器人只需要共享很小的数据包：自己的 odometry 位姿和 AP 估计。

\begin{warnbox}{关键假设}
论文假设机器人之间的相对朝向或旋转矩阵可知。作者提到可由初始编队、IMU 磁力计、AoA 或 RSSI bearing 估计获得。但在室内地面机器人上，磁力计容易受金属和建筑结构影响。你的 Go2 项目如果照搬，需要用 yaw 初始化、视觉标志、UWB bearing 或短程对准动作来解决朝向问题。
\end{warnbox}

\section{数据处理与数据流}

HGP-RL 的数据流可以写成：

\begin{lstlisting}
Per robot:
  random walk / exploration
  -> collect (local odom position, WiFi RSSI from common AP)
  -> train sparse GPR RSSI map
  -> hierarchical inference for AP position
  -> share {local odom pose, predicted AP pose}
  -> AP-oriented vector transformation
  -> relative position of neighbor robots
\end{lstlisting}

在仿真实验中，RSSI 由 path-loss 模型生成：

\[
\mathrm{RSSI}({}^{i}p_i)
=
\mathrm{RSSI}_{d0}
-10\eta \log_{10}\|{}^{i}p_i-{}^{i}p_{\mathrm{AP}}\|
-N_s(0,\sigma_s^2),
\]

并额外加入零均值高斯噪声模拟多径。参数包括 \(\mathrm{RSSI}_{d0}=-20\) dBm、路径损耗指数 \(\eta=3\)、shadow fading 标准差 2 dBm、多径噪声标准差 2 dBm。

\section{实验与结果}

论文在 Robotarium 仿真中部署 3 个 unicycle robot，进行 300 次迭代、10 次 trial。比较对象包括标准 GPR、MEGPR、TRN 和 DGORL。

主要结果：

\begin{itemize}
  \item HGP-RL 的 AP 定位误差接近 dense GPR，但推理时间更低。
  \item 相比 sparse GPR/MEGPR，AP 定位误差降低约 36\%。
  \item 相比 MEGPR-Dense，误差降低约 8\%。
  \item HGP-RL 训练时间比 dense GPR/MEGPR 少约 40\%。
  \item 相对定位 RMSE 在高噪声下比 TRN 和 DGORL 更低，最高约 2 倍提升。
  \item 3 到 10 台机器人规模扩展时，RMSE 从 0.073 m 增至 0.124 m，处理时间从 33 ms 增至 121 ms。
  \item ROS 真机实验中，3 台 Turtlebot2e 在 \(10\text{ m}\times13\text{ m}\) 多房间实验室实现 rendezvous，AP 定位精度 \(0.57\pm0.12\) m，相对轨迹误差 \(0.42\pm0.09\) m。
\end{itemize}

\section{对 Go2 项目的可借鉴方法}

这篇最适合转化为本项目的 RSSI baseline：

\begin{enumerate}
  \item 把充电桩、路由器或 BLE beacon 看成 AP/source。
  \item Go2 在采集阶段记录 \((x,y,s,\mathrm{RSSI},\mathrm{IMU})\)。
  \item 先用 SLAM/Nav2/odom 给 RSSI 打标签，训练 GP radio map。
  \item 在线阶段用分层搜索估计 AP 或路径进度 \(s\)。
  \item 若做多机器人或回充，广播 \(\{\mathrm{odom}, \mathrm{AP/dock\ estimate}\}\)，降低通信量。
\end{enumerate}

如果当前阶段只做一维路径进度 \(s\)，可以简化为：

\begin{lstlisting}
Level 1: every 1.0 m along path
Level 2: around best s, every 0.25 m
Level 3: around best s, every 0.05 m
Output: s_hat, uncertainty
\end{lstlisting}

这样比 dense fingerprint 全表搜索更快，也比纯神经网络更可解释。

\chapter{论文二：Robust Indoor Localization with Ranging-IMU Fusion}

\section{基本信息}

\begin{itemize}
  \item 论文：\emph{Robust Indoor Localization with Ranging-IMU Fusion}
  \item 文件：\file{code/paper/2309.08803v1.pdf}
  \item 主题：UWB/ToF 测距、IMU 预积分、NLOS、多径、鲁棒因子图、iSAM2、D-iSAM2。
\end{itemize}

\section{研究背景}

室内定位常用视觉惯性里程计或视觉 SLAM，但相机带来隐私、算力、能耗和光照问题。UWB 测距功耗低、精度高，适合穿戴设备和机器人。但室内无线测距最大问题是 NLOS 和多径：信号绕墙、反射、遮挡后，测得距离往往比真实直线距离更长，而且误差不是普通对称高斯噪声。

许多系统用 EKF 加异常值剔除处理 range + IMU，但过滤方法只利用当前和过去的局部信息，一旦长时间没有可靠 LOS 测距，IMU 漂移会导致估计偏离，后续正常测距也难以恢复。因此论文选择 smoothing/factor graph，保留较长时间窗口，允许新测量修正过去状态。

\section{核心观点：NLOS 测距误差是单边长尾}

这篇论文最重要的建模洞察是：

\begin{center}
\textbf{真实距离是最短路径，NLOS 测距通常只会把距离测长，而不是对称地测短。}
\end{center}

因此使用普通 Gaussian、Huber 或对称 Cauchy 都不够贴合物理过程。论文把 range measurement likelihood 建模成两种情况的混合：

\begin{itemize}
  \item LOS：测距误差近似高斯，标准差约 0.1--0.2 m。
  \item NLOS：误差为非负、重尾，论文使用 half-Cauchy 分布。
\end{itemize}

NLOS half-Cauchy 的形式可写作：

\[
p(r'-\bar r\mid m=1)\sim
\frac{2}{\pi\gamma\left(1+\left((r'-\bar r)/\gamma\right)^2\right)},
\quad r'-\bar r\ge 0.
\]

其中 \(r'\) 是测量距离，\(\bar r\) 是真实距离，\(\gamma\) 是尺度参数。直觉上，小 NLOS 误差常见，大 NLOS 误差也可能发生，且方向主要是正偏差。

\section{IMU 模型}

论文使用标准 strapdown IMU 模型：

\[
\dot R_t=R_t[\omega_t]_\times,\quad
\dot V_t=g+R_t a_t,\quad
\dot X_t=V_t.
\]

IMU bias 包括加速度计 bias 和陀螺仪 bias：

\[
b_t=[b_t^a,b_t^g]^T.
\]

bias 被建模为一阶 Gauss-Markov 随机过程，并在因子图中通过 bias constraint factor 连接相邻时刻。

IMU 不直接逐点加入优化，而是使用 IMU preintegration，在两个 range timestamp 之间把高频 IMU 累积成一个预积分因子。这是视觉惯性/惯性导航中非常常见的做法，可以避免每个 IMU 样本都引入一个状态节点。

\section{因子图模型}

状态包括：

\begin{itemize}
  \item 每个 range 时刻的设备位姿 \(R_i,X_i\)；
  \item 速度 \(V_i\)；
  \item IMU bias \(b_i\)；
  \item UWB anchor 位置 \(A_j\)，这些位置有先验。
\end{itemize}

因子包括：

\begin{itemize}
  \item prior factor：初始状态与 anchor 先验；
  \item IMU preintegration factor：连接相邻状态；
  \item IMU bias factor：约束 bias 漂移；
  \item range factor：设备到 anchor 的距离观测，使用非对称鲁棒噪声模型。
\end{itemize}

数据流如下：

\begin{lstlisting}
High-rate IMU (~800 Hz)
  -> preintegration between range timestamps

UWB SS-TWR range (1-40 Hz)
  -> range factor to anchors

Initial 10 s:
  -> batch factor graph optimization
  -> initialize iSAM2

Online:
  -> add preintegration and range factors
  -> D-iSAM2 incremental optimization
  -> pose estimates at UWB rate
  -> IMU extrapolation for high-rate pose
\end{lstlisting}

\section{D-iSAM2：给 iSAM2 加简化 trust region}

iSAM2 依赖类似 Gauss-Newton 的线性化更新，在强非线性或条件数差时会不稳定。论文提出 D-iSAM2，本质是在 iSAM2 更新里加入类似 Levenberg-Marquardt 的 damping：

\[
(J^TJ+\lambda I)\delta=J^Tz.
\]

可以理解为给变量加一个特殊因子，误差为 0，但 Jacobian 为 \(\sqrt{\lambda}I\)。若非线性误差下降，则减小 \(\lambda\)；若不下降，则增大 \(\lambda\) 并重试。论文只调整最新变量的 trust region 半径，避免全局重分解 Bayes tree，保持增量效率。

\section{实验与结果}

实验硬件包括 Project Aria 设备上的 IMU 和 DWM3000 UWB 模块，IMU 频率约 800 Hz，UWB 最多 4 个 anchor，每个 anchor 10 Hz，总 range 最高 40 Hz。实验环境为 \(30\text{ m}\times50\text{ m}\) 真实办公室，7 次轨迹，每次 1--3 分钟、50--130 m。

主要结果：

\begin{itemize}
  \item 40 Hz range 下，Proposed + D-iSAM2 多数 run 的平均/最大 3D APE 优于 Gaussian、Huber、Cauchy 等基线。
  \item 在 1 Hz range 下，Huber 多数序列发散；Gaussian 不发散但误差大；Proposed 仍能保持相对稳定。
  \item D-iSAM2 平均更新耗时约 19.21 ms，接近 vanilla iSAM2 的 21.33 ms，显著快于 RISE 的 66.85 ms。
  \item 论文宣称真实办公室平均绝对精度约 0.3 m，并且 range 降到 1 Hz 仍可保持可用精度。
\end{itemize}

\section{对 Go2 项目的可借鉴方法}

本项目当前以 RSSI/IMU 为主，但这篇论文提供了很重要的建模思想：

\begin{enumerate}
  \item \textbf{无线测量误差不要默认高斯。} RSSI、UWB、WiFi FTM、BLE 都会有遮挡和多径，异常值往往具有偏置和长尾。
  \item \textbf{IMU 是低频无线测量之间的桥。} RSSI 或 WiFi scan 可能只有 1 Hz，Go2 IMU/odom 可提供高频运动先验。
  \item \textbf{在线估计要允许回看过去。} 如果只做 EKF，错误 RSSI 峰值可能把状态拉偏；固定滞后平滑或因子图能用后续观测修正过去。
  \item \textbf{异常值模型可从 half-Cauchy 思路迁移。} 对 RSSI 来说异常方向不一定是“距离测长”，但可以建模为 heavy-tailed likelihood 或 switchable factor。
\end{enumerate}

一个适合 Go2 的简化因子图可写成：

\[
x_k=[s_k,\dot s_k,b_k],
\]

其中 \(s_k\) 是路径进度，\(\dot s_k\) 是沿路径速度，\(b_k\) 是 RSSI bias 或设备偏置。因子包括：

\begin{itemize}
  \item odom/IMU motion factor：约束 \(s_{k+1}-s_k\)；
  \item RSSI fingerprint factor：约束当前 RSSI 与 radio map 的匹配；
  \item robust factor：对多径/遮挡 RSSI 使用长尾核；
  \item monotonic factor：若实验为单向路径，约束 \(s\) 不后退或限制最大跳变。
\end{itemize}

\chapter{论文三：SLAM 辅助 WiFi RSSI 指纹数据集构建与 DNN 定位}

\section{基本信息}

\begin{itemize}
  \item 论文：\emph{An Adaptive Indoor Localization Approach Using WiFi RSSI Fingerprinting with SLAM-Enabled Robotic Platform and Deep Neural Networks}
  \item 文件：\file{code/paper/2407.09242v2.pdf}
  \item 主题：WiFi 指纹、SLAM Toolbox、ROS2、机器人自动采集、DTW 时间对齐、DNN 室内定位。
\end{itemize}

\section{研究背景}

WiFi 指纹定位通常分成两个阶段：

\begin{itemize}
  \item offline：在环境中采集不同位置的 WiFi RSSI，构建 fingerprint map；
  \item online：未知位置设备扫描 WiFi，把 RSSI 输入模型，估计位置。
\end{itemize}

瓶颈在 offline 阶段。传统方法要人工在网格点逐点采样并记录位置，耗时、低密度、难更新。环境变化后，指纹数据集失效，重新采集代价高。论文提出用带 LiDAR/IMU/轮速计的 ROS2 机器人，通过 SLAM 自动获得机器人位姿，再把 WiFi 扫描与 SLAM 位姿对齐，自动生成稠密 RSSI 指纹数据集。

\section{系统模型}

系统由四部分组成：

\begin{enumerate}
  \item 商用移动机器人平台：Yahboom Rosmaster X3，Jetson Orin Nano，2D LiDAR，Intel 8265 WiFi，IMU，编码器，RGB-D 相机。
  \item ROS2 数据采集程序：Python controller + FastAPI web interface，控制 survey 开始/停止，记录 ROS2 topics。
  \item SLAM Toolbox：订阅 odom 与 LiDAR scan，输出 map、tf 和机器人在 map 中的位姿。
  \item DNN 定位模型：输入所有 AP 的 RSSI 向量，输出二维位置 \((x,y)\)。
\end{enumerate}

\section{数据处理}

这篇论文最值得本项目借鉴的是数据处理流程。

\subsection{WiFi 扫描}

WiFi 扫描频率约 1 Hz。论文为了减少扫描间隔，选择 1、6、11 这几个非重叠且常用的 WiFi 信道，减少无效扫描时间。

每个 WiFi scan 包含多个 AP 的 RSSI。最终 CSV 格式为：

\begin{lstlisting}
Timestamp, X_Pos, Y_Pos, MAC1, MAC2, ..., MACN
1707935831.6001, 0.0000, 0.0000, 66.0, ..., NaN
1707935832.5993, 0.0026, 0.0034, 66.0, ..., NaN
\end{lstlisting}

注意论文中 RSSI 看起来以正数存储，训练前对缺失 AP 填 \(-100\) dB，表示极弱信号。

\subsection{SLAM 位姿}

机器人通过 wheel encoder + IMU 给出 odom，SLAM Toolbox 用 2D LiDAR scan matching 修正 odom 漂移，并输出 map frame 下的机器人位姿。这个位姿作为 WiFi 指纹的标签。

\subsection{DTW 时间对齐}

WiFi scan 和位姿数据频率不同：WiFi 约 1 Hz，pose/odom 可达 100 Hz，而且开始时间不同。论文使用 Dynamic Time Warping (DTW) 对齐两条时间序列。DTW 通过动态规划寻找最小累计距离路径，允许时间轴拉伸或压缩。

从工程上看，DTW 的作用是：

\begin{itemize}
  \item 给每个 WiFi scan 找到最合理的机器人位姿；
  \item 比简单最近时间戳匹配更能处理采样不均匀或启动延迟；
  \item 但在实时系统中，若时间戳准确，近邻匹配或插值通常更简单。
\end{itemize}

\section{DNN 模型}

模型输入维度等于实验中检测到的 AP 数量。输出维度为 2，对应 map 中的 \(x,y\)。网络结构：

\begin{itemize}
  \item layer sizes：256, 128, 32, 2；
  \item activation：前三层 ReLU，输出层 Linear；
  \item optimizer：Adam；
  \item mini-batch size：32；
  \item early stopping patience：5；
  \item 损失函数文字中提到 MAE，表格中写 MSE，论文这里存在不一致。
\end{itemize}

模型训练后，在线阶段只需要扫描当前 WiFi，按固定 AP 顺序组成向量，缺失 AP 填弱信号值，输入 DNN 得到位置估计。

\section{实验与结果}

实验环境是约 \(4\text{ m}\times10\text{ m}\) 的办公室，面积 40 平方米，检测到 38 个 AP。论文构建了两种人工 ground truth 网格：0.99 m 和 0.66 m 间距，用来验证机器人自动采集结果。

主要结果：

\begin{itemize}
  \item 自动方法采集 320 个 reference points，耗时 320 s，约 1 RP/s，密度 7.27 RP/m\(^2\)。
  \item baseline coarse 采集 23 个点，耗时 656 s，约 0.035 RP/s。
  \item baseline dense 采集 47 个点，耗时 891 s，约 0.052 RP/s。
  \item 自动方法比传统方法构建超过 6 倍密度的数据集，时间效率约 19 倍。
  \item 定位精度 RMSE 0.27 m，MAE 0.19 m，优于表中对比方法。
  \item 使用一半训练数据时误差增加约 60\%，使用四分之一时误差增加约 160\%，说明指纹点密度非常关键。
\end{itemize}

\section{数据流总结}

\begin{lstlisting}
Robot survey:
  ROS2 controller starts recording
  -> LiDAR + encoder + IMU
  -> SLAM Toolbox produces map pose
  -> WiFi scanner records AP RSSI at ~1 Hz
  -> ROS bag / logs
  -> DTW aligns WiFi scans with SLAM poses
  -> CSV fingerprint dataset
  -> DNN training
  -> online WiFi localization
\end{lstlisting}

\section{对 Go2 项目的可借鉴方法}

这篇论文和你的项目最直接相关，因为你已经有 Go2、ROS2、RSSI/IMU 采集和路径实验。建议照搬它的数据集思想，但按你的项目目标做改造：

\begin{enumerate}
  \item 不一定先做二维 \((x,y)\)，可以先做路径进度 \(s\)。
  \item 用 Nav2/SLAM/odom 生成每个 RSSI window 的标签：\((s,x,y,\theta)\)。
  \item CSV 推荐格式：
\end{enumerate}

\begin{lstlisting}
timestamp,s,x,y,yaw,vx,wz,ax,ay,az,gx,gy,gz,
beacon_1_mean,beacon_1_std,...,beacon_N_mean,beacon_N_std
\end{lstlisting}

与当前项目中的 \file{README.md} 和 baseline 格式相比，可以保留已有 \code{beacon_*_mean/std}，但加入 SLAM/odom 生成的 \code{x,y,yaw,s}，这样同一份数据可以训练：

\begin{itemize}
  \item \(RSSI\rightarrow s\) 的路径定位模型；
  \item \(RSSI\rightarrow (x,y)\) 的二维定位模型；
  \item \(RSSI+IMU\rightarrow s\) 的序列模型；
  \item GP/HGP radio map baseline。
\end{itemize}

\chapter{论文四：Radar Teach and Repeat}

\section{基本信息}

\begin{itemize}
  \item 论文：\emph{Radar Teach and Repeat: Architecture and Initial Field Testing}
  \item 文件：\file{code/paper/2409.10491v1.pdf}
  \item 主题：FMCW 扫描雷达、Teach-and-Repeat、radar-gyro odometry、ICP、local submap、MPC、长期导航。
\end{itemize}

\section{研究背景}

移动机器人在巡逻、农业、矿山、校园 shuttle、重复巡检等任务中，经常沿固定路线反复行驶。Teach-and-Repeat (T\&R) 的思想是：先人工开一遍安全路线并建局部地图，之后机器人重复这条路线。相比全局 SLAM，T\&R 不强制构建全局一致大地图，而是用沿路径串联的局部子图，适合长期重复任务。

传统 T\&R 常用视觉或 LiDAR：

\begin{itemize}
  \item 视觉怕光照、黑暗、视角变化；
  \item LiDAR 不怕光照，但容易受雪、雾、尘、烟、植被等影响；
  \item GPS 在树冠、室内、城市峡谷下不可靠。
\end{itemize}

FMCW 雷达波长更长，对尘雾雪烟更鲁棒，且量测距离远，适合恶劣环境长期导航。论文实现了完整 closed-loop Radar Teach-and-Repeat，声称是首次用 radar 作为唯一外感知传感器在 off-road 环境中闭环重复路线。

\section{系统模型}

传感器：

\begin{itemize}
  \item Navtech RAS3 2D scanning radar，4 Hz；
  \item single-axis yaw gyro，100 Hz；
  \item RTK GPS 仅用于评估；
  \item Ouster LiDAR 用作对比基线，不用于 radar repeat。
\end{itemize}

系统分为 teach phase 和 repeat phase。

\section{Teach Phase}

\subsection{雷达点云提取}

雷达原始数据是极坐标 radar image。论文使用改造的 BFAR detector，在每个 azimuth 中提取关键点。检测点经过聚类/质心处理后转为 2D Cartesian point cloud，用于 ICP odometry 和地图构建。

\subsection{Radar-gyro 连续时间 ICP 里程计}

每个雷达 scan 到来时，系统优化一个包含姿态和速度的状态：

\[
x_i=\{T_i,\varpi_i\}.
\]

代价函数包括三类误差：

\[
J_{\mathrm{odom,radar}}
=
\frac{1}{2}e_{\mathrm{prior}}^TQ_{\mathrm{prior}}^{-1}e_{\mathrm{prior}}
+
\sum_{j=1}^{M}\frac{1}{2}e_{\mathrm{odom},j}^TR_j^{-1}e_{\mathrm{odom},j}
+
\frac{1}{2}e_{\mathrm{yaw}}^TR_{\mathrm{yaw}}^{-1}e_{\mathrm{yaw}}.
\]

其中：

\begin{itemize}
  \item motion prior 假设短时间内近似恒速；
  \item radar ICP error 做当前 scan 与滑动地图/局部点云的点到点匹配；
  \item yaw preintegration factor 聚合两个 radar scan 之间的陀螺仪 yaw rate。
\end{itemize}

这里和 2309 论文有共同点：低频外部观测之间，用高频惯性测量补足运动约束。

\subsection{Mapping}

人工 teach 时，系统把一段窗口内的雷达点云合并成 local submap。当机器人移动距离或转角超过阈值时创建新 submap。相邻 submap 通过相对里程计连接，形成沿路径的 pose graph。这个图不追求全局一致性，而追求局部可重复定位。

\section{Repeat Phase}

\subsection{High-rate gyro odometry}

雷达只有 4 Hz，直接用于控制太慢。论文在两个雷达 scan 之间，用 100 Hz gyro 解一个小优化问题，提供高频 odometry。新雷达 scan 到来后，再用更准的 ICP odometry 替换高频外推结果。

这对 Go2 很重要：RSSI/WiFi scan 甚至可能只有 1 Hz，更需要 IMU/odom 在中间补足连续状态。

\subsection{Localization ICP}

repeat 阶段，当前雷达 scan 与 teach 阶段保存的最近 local submap 做 ICP，求机器人当前 frame 到 teach path local map frame 的相对位姿。代价函数由 pose prior 和点云 ICP 误差组成。

\subsection{MPC 闭环控制}

Localization 结果送入 MPC。MPC 以 teach path 上未来 \(K\) 个目标位姿为参考，优化线速度 \(v_k\) 和角速度 \(\omega_k\)，满足：

\begin{itemize}
  \item unicycle motion model；
  \item 速度上下限；
  \item corridor constraint，即机器人不能偏离 teach path 太远；
  \item 若约束不可行，机器人停下等待人工干预。
\end{itemize}

论文中 corridor 最大 lateral deviation 设置为 50 cm。

\section{实验与结果}

平台为 Clearpath Warthog UGV。四条路线包括 Parking Loop、Mars Dome Loop、Grassy Loop、Woody Loop，总 teach 距离 1276 m，LiDAR repeat 8.1 km，Radar repeat 11.8 km。

主要结果：

\begin{itemize}
  \item RT\&R 在四个环境中 autonomy rate 为 100\%，无人干预完成 11.8 km。
  \item Parking Loop radar lateral RMSE 5.6 cm，最大误差 21.7 cm。
  \item Mars Dome radar lateral RMSE 7.5 cm，最大误差 24.0 cm。
  \item Grassy Loop radar lateral RMSE 12.1 cm，最大误差 43.8 cm。
  \item LiDAR T\&R 在所有路线更准，但 radar 在结构较好的区域接近 LiDAR。
  \item Woody Loop 树冠导致 GPS ground truth 不可靠，无法计算 GPS-measured RMSE，但内部估计仍显示可重复行驶。
\end{itemize}

\section{局限与经验}

论文讨论了 radar 的工程难点：

\begin{itemize}
  \item 雷达点云会出现 radial artifacts，也就是以传感器为中心的稀疏环形假点。
  \item 场景几何稀疏时，ICP 可能错误对齐这些假点，导致 yaw drift。
  \item yaw 误差会和控制器形成坏反馈，机器人可能逐渐偏离路径。
  \item 加入 gyro yaw rate 对抑制 heading drift 很关键。
  \item 未来需要更好的点特征提取和 Doppler correction。
\end{itemize}

\section{对 Go2 项目的可借鉴方法}

如果你的 Go2 要做固定路线 RSSI 采集、巡检或自动回充，T\&R 思想非常适合：

\begin{lstlisting}
Teach:
  human drives Go2 along safe route
  -> record odom/IMU/RSSI/pointcloud/camera
  -> build route map and path coordinate s

Repeat:
  localize to taught route
  -> estimate path progress s and lateral error e
  -> waypoint/Nav2/MPC controls Go2
  -> RSSI model corrects drift or detects route segment
\end{lstlisting}

短期不用上雷达，也可以借鉴架构：

\begin{itemize}
  \item 用 LiDAR/point cloud 或视觉标签替代 radar submap；
  \item 用 Go2 IMU/odom 替代 gyro-only high-rate odometry；
  \item 用 RSSI fingerprint 判断当前处于路线哪一段；
  \item 用 waypoint follower 或 Nav2 controller 做 repeat 控制。
\end{itemize}

\chapter{结论}

四篇论文给本项目的核心启发可以压缩成一句话：

\begin{center}
\textbf{先用 SLAM/Nav2/odom 把 RSSI 数据采好、标准化；再用 RSSI 模型提供低成本位置线索；最后用 IMU/odom/路径约束把低频、噪声大的无线观测融合成连续可控的机器人状态。}
\end{center}

最推荐的近期实现路径是：

\begin{enumerate}
  \item 按 2407 论文思路，完成 Go2 自动/半自动 RSSI 指纹数据采集与标签生成。
  \item 在现有 KNN/DNN baseline 外，增加 2307 式一维 HGP 路径进度搜索。
  \item 用 2309 的鲁棒观测思想，为 RSSI 异常值加入 heavy-tailed loss 或跳变约束。
  \item 若目标扩展到自动回充或固定巡检，借鉴 2409 的 Teach-and-Repeat，将路线表示为 \(s,e,\theta_e\) 并接入 waypoint/Nav2 控制。
\end{enumerate}

\end{document}
